% =================================================================
\chapter{Customizing}
% =================================================================

\section{Disclaimer}
Die Software von LogObject nutzt standardmäßig die hier dokumentierte Logging-Konfiguration.
Diese Standardisierung gewährleistet über alle Kunden hinweg eine konsistente Basis, die schnellen, zuverlässigen und effektiven Support ermöglicht.

Unabgestimmte Änderungen an dieser Konfiguration sind ausdrücklich nicht erwünscht, da sie die Nachvollziehbarkeit und Supportfähigkeit erheblich beeinträchtigen können.

Insbesondere in kritischen Situationen können kundenspezifische Anpassungen der Logging-Kon\-fi\-gu\-ra\-tion zu Verzögerungen oder zusätzlichen Problemen führen.
Daher wird dringend davon abgeraten, die Verzeichnisstrukturen der Logdateien eigenständig zu verändern.

Anpassungen von Logdatei-Mengen, Archivierungsstrategien oder Rollover-Größen sollten ausschließlich nach vorheriger Abstimmung und mit sorgfältiger Planung unter Berücksichtigung der lokalen Gegebenheiten erfolgen. Änderungen ohne Abstimmung erfolgen auf eigenes Risiko und können den Support einschränken.

\section{Weiterführende Informationen}

Für detaillierte Informationen zur Bedienung und Konfiguration des Loggings wird auf die offizielle \texttt{Log4j}-Dokumentation verwiesen.
Dort finden sich Hinweise zu Log-Leveln, Appendern, Layouts, Rollover-Strategien und weiteren Konfigurationsmöglichkeiten.

\section{Löschen leerer Rollover-Ordner}
\label{sec:customizing_cleanup}
Das Löschen leerer Datumsordner (\texttt{yyyy-mm-dd}) ist eine direkte Folge der
kundenspezifisch angepassten Logging-Konfiguration.

Im Rahmen des Customizings wurde die Log4j2-Konfiguration so eingerichtet, dass Logs
täglich in datumsbezogene Unterordner archiviert werden. Das löschen der alten Logs sorgt dafür das die Datumsordner leer auf dem Server liegen.

Hierfür  ist das Script \texttt{deleteEmptyFolders.sh} in Kombination mit dem in Abschnitt~\ref{sec:cleanup} beschriebenen Cron-Job.
